This chapter introduces the research context, motivation, objectives, and method. It
also describes the scope and main limitations of the study.

\section{Background}
Modern observatories collect high-cadence data at several wavelengths. Missions such
as \textit{Chandra}, \textit{XMM-Newton}, and \textit{NICER} have greatly increased the
amount and quality of available astrophysical data \cite{weisskopf2002, jansen2001,
gendreau2012}. These data reveal complex changes in black hole accretion systems.

Black holes are regions of spacetime from which light cannot escape after crossing
the event horizon. The matter around them can still produce strong radiation.
Accretion occurs when matter is captured and spirals inward. It converts gravitational
energy into heat and radiation with high efficiency \cite{shakura1973, frank2002}.
Black hole binaries and Active Galactic Nuclei (AGN) vary over several timescales
\cite{remillard2006, done2007}. Their variability reflects changes in the geometry of
the flow and in the mechanism that releases energy. These changes are commonly
described using the \textit{Hard State} and \textit{Soft State}
\cite{fender2004}.

Accretion systems also follow power-law scaling relations across a wide range of
black-hole masses. These relations suggest shared physical behaviour and motivate
the Fundamental Plane of black hole activity \cite{merloni2003}. However, correlation
alone cannot show which physical process drives another. Causal discovery provides
a way to study these directed relationships.

The main difficulty is that the physical drivers of accretion, such as magnetic-field
configuration, viscosity, and local shocks, are usually not observed directly. We
observe only their emitted light. This thesis studies \textbf{Latent Regime Discovery}:
finding the hidden physical states and their causal mechanisms from observational
data.

\section{Rationale of the Study}
Black hole accretion systems change between physically distinct states, but the state
of each observation is not labelled. Manual classification is time-consuming and does
not scale well to high-cadence instruments such as \textit{NICER}.

The causal relationships can also change between states. A pooled model may therefore
produce biased coefficients and hide the true jet-disk coupling. For example, the
relationship between X-ray luminosity and radio emission may differ between the Hard
State and the Soft State. That difference cannot be recovered unless the observations
are first separated by regime.

An automated framework for latent causal regime discovery would reduce the need for
manual labelling and could reveal physical relationships that correlation-based
analyses miss.

\section{Problem Statement}
This thesis addresses the lack of methods that can discover latent causal regimes
from unlabelled observations without environment labels, a surrogate domain index,
or a pre-specified number of states. ICP \cite{peters2016causal} and IRM
\cite{arjovsky2019irm} require environment labels. CD-NOD \cite{huang2020causal}
requires an auxiliary time or domain index. CCSL requires multiple independent
subject datasets, whereas X-ray binary monitoring usually provides one continuous
record.

Traditional clustering groups observations by feature similarity rather than causal
mechanism. It can miss boundaries where observed distributions overlap but the causal
laws differ. High-cadence accretion data also contain temporal autocorrelation.
Methods that assume independent and identically distributed (i.i.d.) observations
may then produce inflated statistics and false regime boundaries.

The problem is therefore the joint treatment of unlabelled data, an unknown regime
count, temporal dependence, and invariance-based recursive partitioning.
\section{Objectives}
The general objective is to develop and validate a framework that discovers hidden
causal regimes and their structures from unlabelled, temporally dependent data. The
framework should not require environment labels, a surrogate nonstationarity index,
or a pre-specified number of states. The specific objectives are:
\begin{enumerate}
    \item To formulate accretion observables within a structural causal model grounded
    in established scaling physics, with parameters that have a physical interpretation.
    \item To design a statistical criterion that distinguishes genuine mechanism
    changes from variation caused by noise.
    \item To maintain the validity of this criterion when high-cadence observations
    are temporally dependent.
    \item To determine the number and location of latent regimes automatically.
    \item To prevent overfitting and unnecessary fragmentation of the time series.
    \item To recover a directed causal structure within each regime and compare the
    coupling between physical variables across states.
    \item To validate the framework on data with known ground-truth regimes and graphs
    before applying it to real observations.
\end{enumerate}

\section{Methodology in Brief}
The method has two stages. First, the observables are transformed into log space so
that the power-law relationships become linear. A global OLS model is then fitted,
and a Chow $F$-test \cite{chow1960} checks whether one causal mechanism explains the
series. A first-order autoregressive term accounts for temporal autocorrelation,
giving a lagged Chow statistic with $q=6$ parameters per model.

If invariance is rejected, the method searches for the strongest candidate split.
The Bayesian Information Criterion (BIC) \cite{schwarz1978} decides whether the
split improves the model enough to justify the added complexity. This process is
applied recursively, so the number of regimes $K$ is determined automatically.
Second, DYNOTEARS estimates a directed causal graph within each stable partition.
The framework is first tested on synthetic structural causal models with known
graphs. Regime recovery is evaluated with the Adjusted Rand Index
\cite{hubert1985}. Future work will apply it to the WATCHDOG catalogue
\cite{tetarenko2016} and the Fundamental Plane compilation of Merloni et al.
\cite{merloni2003}.

The overall methodology is summarised in Figure~\ref{fig:pipeline_intro}.

\begin{figure}[h]
    \centering
    \begin{tikzpicture}[
        node distance=0.7cm,
        every node/.style={align=center, font=\small},
        outerbox/.style={draw=black!60, rectangle, rounded corners=1ex, inner sep=10pt, text width=8.5cm, fill opacity=0.9},
        arrow/.style={->, >=Stealth, thick, draw=black!80}
    ]

    % 1. Data Node
    \node[outerbox, fill=cyan!30] (data) {\textbf{Unlabeled Observational Data} \\ \vspace{0.3em} \footnotesize Multiwavelength accretion time-series data};

    % 2. Prep Node
    \node[outerbox, fill=teal!30] (prep) [below=of data] {\textbf{Data Preprocessing} \\ \vspace{0.3em} \footnotesize Log-space transformation \& temporal lag integration};

    % 3. Stage 1 Node
    \node[outerbox, fill=blue!20] (stage1) [below=of prep] {\textbf{Stage 1: Latent Regime Discovery} \\ \vspace{0.3em} \footnotesize Recursive timeline splitting via Binary Segmentation \\ \footnotesize Evaluated using the Lagged Chow F-Test \& BIC Penalty};

    % 4. Stage 2 Node
    \node[outerbox, fill=orange!30] (stage2) [below=of stage1] {\textbf{Stage 2: Causal Graph Recovery} \\ \vspace{0.3em} \footnotesize Extracting localized Window Causal Graphs using DYNOTEARS};

    % 5. Output Node
    \node[outerbox, fill=purple!30] (output) [below=of stage2] {\textbf{Output: K-Regime-Specific Causal Graphs} \\ \vspace{0.3em} \footnotesize Dynamic causal rules mapping distinct physical accretion states};

    % Main Vertical Arrows
    \draw[arrow] (data) -- (prep);
    \draw[arrow] (prep) -- (stage1);
    \draw[arrow] (stage1) -- (stage2);
    \draw[arrow] (stage2) -- (output);

    \end{tikzpicture}
    \caption{High-level pipeline of the unsupervised dynamic causal discovery framework, illustrating the progression from observational data to regime-specific causal graphs.}
    \label{fig:pipeline_intro}
\end{figure}
\vspace{1.5em}

\section{Scopes and Challenges}
This study combines causal inference with high-energy observational astrophysics.
The framework is designed for multi-wavelength black-hole accretion data and uses
log-linear structural equations within each regime. Its outputs are regime-specific
causal graphs, regression coefficients, and autoregressive parameters.

Several limitations remain. The linear model may miss nonlinear relationships, and
the piecewise-stationary model is only an approximation to continuous accretion
dynamics. The Chow test and BIC derivations assume Gaussian noise. The temporal
model uses only one lag, and the current validation uses synthetic rather than real
observations. Unmeasured variables, such as magnetic-field geometry, may create
spurious associations that the PC algorithm cannot resolve without additional
knowledge \cite{spirtes2000causation}. Residual autocorrelation may also inflate the
Chow statistic; Newey-West standard errors \cite{newey1987} are considered as a
robustness check. Non-Gaussian noise and violations of faithfulness provide further
challenges for causal discovery.
